% ═══════════════════════════════════════════════════════════════
% MUSTERLÖSUNG: Las preposiciones de lugar
% Spanisch Klasse 8 - Unidad 6, DS 2 (Einzelstunde)
% ═══════════════════════════════════════════════════════════════

\documentclass[10pt, a4paper]{article}

\usepackage[utf8]{inputenc}
\usepackage[T1]{fontenc}
\usepackage[ngerman]{babel}

% --- SCHRIFTEN ---
\usepackage{palatino}
\usepackage{helvet}
\newcommand{\sanstext}[1]{{\fontfamily{phv}\selectfont #1}}

\usepackage{microtype}
\usepackage[margin=1.4cm, top=1.6cm, bottom=1.3cm]{geometry}
\usepackage{enumitem}
\usepackage{tabularx}
\usepackage{booktabs}
\usepackage{array}
\usepackage{multicol}
\usepackage{tikz}
\usetikzlibrary{calc}

\usepackage{fancyhdr}
\usepackage{lastpage}
\usepackage{needspace}

\usepackage{xcolor}
\usepackage{tcolorbox}
\tcbuselibrary{skins}

% ═══════════════════════════════════════════════════════════════
% FARBEN
% ═══════════════════════════════════════════════════════════════

\definecolor{kopfzeile}{HTML}{1A1A1A}
\definecolor{akzent}{HTML}{C41E3A}
\definecolor{hellgrau}{HTML}{F2F2F2}
\definecolor{mittelgrau}{HTML}{777777}
\definecolor{dunkelgrau}{HTML}{444444}
\definecolor{loesung}{HTML}{C62828}

% ═══════════════════════════════════════════════════════════════
% VARIABLEN
% ═══════════════════════════════════════════════════════════════

\newcommand{\thema}{Las preposiciones de lugar}
\newcommand{\klasse}{8}
\newcommand{\maxpunkte}{20}

% ═══════════════════════════════════════════════════════════════
% KOPF- UND FUSSZEILE
% ═══════════════════════════════════════════════════════════════

\pagestyle{fancy}
\fancyhf{}
\renewcommand{\headrulewidth}{0pt}
\renewcommand{\footrulewidth}{0pt}
\cfoot{\sanstext{\footnotesize\textcolor{mittelgrau}{\thepage\,/\,\pageref{LastPage}}}}

% ═══════════════════════════════════════════════════════════════
% BOXEN
% ═══════════════════════════════════════════════════════════════

\newtcolorbox{vokabelkasten}{
    colback=white,
    colframe=mittelgrau,
    boxrule=0.5pt,
    arc=0pt,
    left=8pt, right=8pt, top=6pt, bottom=6pt
}

\newtcolorbox{hinweisbox}{
    colback=hellgrau,
    colframe=hellgrau,
    boxrule=0pt,
    arc=0pt,
    left=8pt, right=8pt, top=6pt, bottom=6pt,
    borderline north={1.5pt}{0pt}{dunkelgrau}
}

% ═══════════════════════════════════════════════════════════════
% BEFEHLE
% ═══════════════════════════════════════════════════════════════

\newcommand{\es}[1]{\textbf{#1}}
\newcommand{\de}[1]{{\small\textit{\textcolor{mittelgrau}{#1}}}}
\newcommand{\lsg}[1]{\textcolor{loesung}{\textbf{#1}}}
\newcommand{\punktebox}[1]{\hfill\sanstext{\small[\_\_/#1]}}

\newcommand{\aufgabe}[2]{%
    \needspace{4\baselineskip}%
    \vspace{12pt}%
    \noindent\sanstext{\textbf{#1}\quad#2}%
    \vspace{6pt}%
    \nopagebreak[4]%
}

\setlength{\parindent}{0pt}
\setlength{\parskip}{4pt}
\setlength{\columnsep}{24pt}

% ═══════════════════════════════════════════════════════════════
% DOKUMENT
% ═══════════════════════════════════════════════════════════════

\begin{document}

% ═══════════════════════════════════════════════════════════════
% HEADER
% ═══════════════════════════════════════════════════════════════

\noindent
\begin{tikzpicture}[remember picture, overlay]
    \fill[kopfzeile] (current page.north west) rectangle +(\paperwidth, -1cm);
    \node[anchor=west, text=white, font=\sanstext\large\bfseries]
        at ([xshift=1.4cm, yshift=-0.5cm]current page.north west) {\thema{} -- \textcolor{loesung}{MUSTERLÖSUNG}};
    \node[anchor=east, text=white, font=\sanstext\small]
        at ([xshift=-1.4cm, yshift=-0.5cm]current page.north east)
        {Spanisch\,·\,Klasse \klasse\,·\,ML};
\end{tikzpicture}

\vspace{0.5cm}

% ═══════════════════════════════════════════════════════════════
% PRÄPOSITIONENTABELLE (mit Lösungen)
% ═══════════════════════════════════════════════════════════════

\begin{vokabelkasten}
\textbf{Die Ortspräpositionen}

\vspace{6pt}

\begin{multicols}{2}
\renewcommand{\arraystretch}{1.3}
\begin{tabular}{@{}p{4cm}p{3.5cm}@{}}
\es{enfrente de} & \lsg{gegenüber von} \\
\es{encima de} & \lsg{über, auf} \\
\es{al lado de} & \lsg{neben} \\
\es{a la izquierda de} & \lsg{links von} \\
\es{en} & \lsg{in, an, auf} \\
\end{tabular}

\columnbreak

\begin{tabular}{@{}p{4cm}p{3.5cm}@{}}
\es{a la derecha de} & \lsg{rechts von} \\
\es{entre} & \lsg{zwischen} \\
\es{detrás de} & \lsg{hinter} \\
\es{debajo de} & \lsg{unter} \\
\es{delante de} & \lsg{vor} \\
\end{tabular}
\end{multicols}
\end{vokabelkasten}

\vspace{6pt}

% ═══════════════════════════════════════════════════════════════
% AUFGABE 1: Copito
% ═══════════════════════════════════════════════════════════════

\aufgabe{1}{Copito: Beschreibe die Position des Hundes (Buch S. 106, Übung 2B).} \punktebox{10}

{\footnotesize\textit{Hinweis: Die genauen Lösungen hängen von den Buchbildern ab. Beispielantworten:}}

\vspace{4pt}

\begin{multicols}{2}
\begin{tabular}{@{}rl@{}}
Bild 1: & \lsg{Copito está debajo de la mesa.} \\[0.6em]
Bild 2: & \lsg{Copito está encima de la silla.} \\[0.6em]
Bild 3: & \lsg{Copito está delante de la puerta.} \\[0.6em]
Bild 4: & \lsg{Copito está detrás del sofá.} \\[0.6em]
Bild 5: & \lsg{Copito está al lado de la cama.} \\
\end{tabular}

\columnbreak

\begin{tabular}{@{}rl@{}}
Bild 6: & \lsg{Copito está entre las sillas.} \\[0.6em]
Bild 7: & \lsg{Copito está a la derecha de la mesa.} \\[0.6em]
Bild 8: & \lsg{Copito está a la izquierda de la ventana.} \\[0.6em]
Bild 9: & \lsg{Copito está enfrente del televisor.} \\[0.6em]
Bild 10: & \lsg{Copito está en la caja.} \\
\end{tabular}
\end{multicols}

\vspace{4pt}
{\footnotesize\textit{Bewertung:} 1P pro korrektem Satz mit passender Präposition}

\vspace{6pt}

% ═══════════════════════════════════════════════════════════════
% AUFGABE 2: Barrio beschreiben
% ═══════════════════════════════════════════════════════════════

\aufgabe{2}{Mi barrio: Beschreibe ein Stadtviertel (Buch S. 107, Übung 4).} \punktebox{10}

\begin{hinweisbox}
\textbf{Schritt 1:} Audio-Lösung hängt vom Hörtext ab.

\vspace{4pt}

\textbf{Schritt 2: Mögliche Sätze für Barrio-Beschreibung:}
\end{hinweisbox}

\vspace{6pt}

\begin{enumerate}[label=\arabic*., itemsep=4pt]
    \item \lsg{La iglesia está enfrente de la plaza.}
    \item \lsg{El supermercado está al lado de la farmacia.}
    \item \lsg{La biblioteca está a la derecha del parque.}
    \item \lsg{El cine está entre el restaurante y la cafetería.}
    \item \lsg{La panadería está a la izquierda del banco.}
    \item \lsg{(Zusatz) El teatro está detrás del ayuntamiento.}
\end{enumerate}

\vspace{6pt}

{\footnotesize\textit{Bewertung:} 2P pro korrektem Satz (1P Struktur, 1P passende Präposition). Mindestens 5 Sätze erforderlich.}

% ═══════════════════════════════════════════════════════════════
% FOOTER
% ═══════════════════════════════════════════════════════════════

\vfill

\noindent\rule{\textwidth}{0.5pt}

\vspace{4pt}

\hfill\sanstext{\textbf{Maximalpunktzahl:}} \maxpunkte

\end{document}
